\documentclass{article}
\usepackage{amsmath, amssymb, amsfonts}
\usepackage{mathrsfs}

% Custom operators and macros
\newcommand{\Pic}{\mathbf{Pic}}
\newcommand{\pPic}{\operatorname{Pic}}
\newcommand{\Hilb}{\mathbf{Hilb}}
\newcommand{\Div}{\operatorname{Div}}
\newcommand{\LinSys}{\operatorname{LinSys}}
\newcommand{\et}{\text{ét}}
\newcommand{\bEGA}{\textbf{EGA }}

\newtheorem{theorem}{THEOREM}[section]

\begin{document}

\begin{theorem}[Main]
Assume $f : X \to S$ is projective Zariski locally over $S$, and is flat with integral geometric fibers.
\begin{enumerate}
    \item Then $\Pic_{X/S}$ exists, is separated and locally of finite type over $S$, and represents $\pPic_{(X/S)(\et)}$.
    \item If also $S$ is Noetherian and $X/S$ is projective, then $\Pic_{X/S}$ is a disjoint union of open subschemes, each an increasing union of open quasi-projective $S$-schemes.
\end{enumerate}
\end{theorem}

\noindent\textbf{PROOF.} By [\textbf{EGA G} (0, 4.5.5), p. 106], it is a local matter on $S$ to represent a Zariski sheaf on the category of $S$-schemes. Moreover, it is also a local matter on $S$ to prove that an $S$-scheme is separated and locally of finite type. Hence, in order to prove (1), we may assume $S$ is Noetherian and $X/S$ is projective.

Plainly an $S$-scheme is separated if it is a disjoint union of separated open subschemes, or if it is an increasing union of separated open subschemes. Hence (1) follows from (2).

To prove (2), owing to Yoneda's lemma, we may view the category of schemes as a full subcategory of the category of functors by identifying a scheme $T$ with its functor of points. Denote this functor too by $T$ in order to lighten the notation. And say that the functor is a scheme, as well as that it is representable. Also, set
\[
P := \pPic_{(X/S)(\et)}.
\]
Note $P(T) = \operatorname{Hom}(T, P)$.

Given a polynomial $\phi \in \mathbb{Q}[n]$, let $P^{\phi} \subset P$ be the \'{e}tale subsheaf associated to the presheaf whose $T$-points are represented by the invertible sheaves $\mathcal{L}$ on $X_T$ such that we have
\begin{equation}\label{4.8.1}
\chi(X_t, \mathcal{L}_t^{-1}(n)) = \phi(n) \text{ for all } t \in T.
\end{equation}
Notice that this presheaf is well defined, because \eqref{4.8.1} remains valid after any base change $p: T' \to T$; indeed, for any $t' \in T'$, for any $i$, and for any $n$, we have
\[
\mathrm{H}^i(X_{t'}, \mathcal{L}_{t'}^{-1}(n)) = \mathrm{H}^i(X_{p(t')}, \mathcal{L}_{p(t')}^{-1}(n)) \otimes_{k_t} k_{t'}
\]
because cohomology commutes with flat base change by [\textbf{Ha83}, Prp. III, 9.3, p. 255]. Hence $P^{\phi} $ is well defined too.

Fix a map $T \to P$, and represent it by means of an \'{e}tale covering $p : T' \to T$ and an invertible sheaf $\mathcal{L}'$ on $X_{T'}$. Consider the subset $T^{\prime\phi} \subset T'$ defined as follows:
\[
T^{\prime\phi} := \{ t' \in T' \mid \chi(X_{t'}, \mathcal{L}_{t'}^{\prime-1}(n)) = \phi(n) \}.
\]
Then $T^{\prime\phi}$ is open by [\textbf{EGA III}$_2$, 7.9.11].

Set $T^{\phi} := p(T^{\prime\phi})$. Then $T^{\phi} \subset T$ is open as $T^{\prime\phi} \subset T'$ is open and $p$ is \'{e}tale. Moreover, $T^{\prime\phi} = p^{-1}(T^{\phi})$. Indeed, let $t' \in p^{-1}(T^{\phi})$. Say $p(t') = p(t'_1)$ where $t'_1 \in T^{\prime\phi}$. Now, there is an \'{e}tale covering $T'' \to T' \times_T T'$ such that the two pullbacks of $\mathcal{L}'$ to $X_{T''}$ are isomorphic. Let $t'' \in T''$ have image $t' \in T'$ under the first map $T'' \to T'$ and have the image $t'_1 \in T'$ under the second map. Then
\[
\chi(X_{t'}, \mathcal{L}_{t'}^{\prime-1}(n)) = \chi(X_{t''}, \mathcal{L}_{t''}^{\prime\prime-1}(n)) = \chi(X_{t'_1}, \mathcal{L}_{t'_1}^{\prime-1}(n)) = \phi(n).
\]
Hence $t' \in T^{\prime\phi}$. Thus $T^{\prime\phi} \supset p^{-1}(T^{\phi})$. Therefore, $T^{\prime\phi} = p^{-1}(T^{\phi})$.

Furthermore, $T^{\phi}$ is (represents) the fiber product of functors $P^{\phi} \times_P T$. Indeed, to see they have the same $R$-points, let $r : R \to T$ be a map; form $R' := R \times_T T'$ and $r' : R' \to T'$. Suppose $r$ factors through $T^{\phi}$. Then $r'$ factors through $T^{\prime\phi}$. So $R' \to T' \to P$ factors through $P^{\phi}$ essentially by definition. Now, $R' \to R$ is an \'{e}tale covering. Hence $R \to T \to P$ factors through $P^{\phi}$ since $P^{\phi}$ is an \'{e}tale sheaf.

Conversely, suppose $R \to T \to P$ factors through $P^{\phi}$. Then $R \to P$ is defined by means of an \'{e}tale covering $R'' \to R$ and an invertible sheaf $\mathcal{L}''$ on $X_{R''}$ such that $\chi(X_u, \mathcal{L}_u^{\prime\prime-1}(n)) = \phi(n)$ for all $u \in R''$. Since both $\mathcal{L}''$ and $\mathcal{L}'$ define $R \to P$, there is an \'{e}tale covering $R''' \to R'' \times_R R'$ such that the pullbacks of $\mathcal{L}''$ and $\mathcal{L}'$ to $X_{R'''}$ are isomorphic. Hence the image of $r' : R' \to T'$ lies in $T^{\prime\phi}$. But the latter is open. Hence $r'$ factors through it. Therefore, $r : R \to T$ factors through $T^{\phi}$. Thus $T^{\phi}$ and $P^{\phi} \times_P T$ have the same $R$-points.

Let $\phi$ vary. Plainly the $T^{\prime\phi}$ are disjoint and cover $T'$. So the $T^{\phi}$ are disjoint and cover $T$. Hence, by a general result [\textbf{EGA G} (0, 4.5.4), p. 103], if the $P^{\phi}$ are (representable by) schemes, then $P$ is their disjoint union. Thus it remains to represent each $P^{\phi}$ by an increasing union of open quasi-projective $S$-schemes.

Fix $\phi$. Given $m \in \mathbb{Z}$, let $P^{\phi}_m \subset P^{\phi}$ be the \'{e}tale subsheaf associated to the presheaf whose $T$-points are represented by the $\mathcal{L}$ on $X_T$ such that, in addition to \eqref{4.8.1}, we have
\begin{equation}\label{4.8.2}
\mathrm{R}^i f_{T*} \mathcal{L}(n) = 0 \text{ for all } i \ge 1 \text{ and } n \ge m.
\end{equation}
Notice that this presheaf is well defined, because \eqref{4.8.2} remains valid after any base change $p: T' \to T$, as is shown next.

First, let's see that \eqref{4.8.2} is equivalent to the following condition:
\begin{equation}\label{4.8.3}
\mathrm{H}^i(X_t, \mathcal{L}_t(n)) = 0 \text{ for all } i \ge 1\text{, all } n \ge m\text{, and all } t \in T.
\end{equation}
Indeed, \eqref{4.8.3} implies \eqref{4.8.2}; in fact, for any given $i, t$ and $n$, if $\mathrm{H}^i(X_t, \mathcal{L}_t(n)) = 0$, then $\mathrm{R}^i f_{T*} (\mathcal{L}(n) \otimes_{\mathcal{O}_T} \mathcal{N})_t = 0$ for all quasi-coherent $\mathcal{N}$ on $T$ by [\textbf{EGA III}$_2$, 7.5.3] or [\textbf{OB72}, Cor. 2.1 p. 68].

Conversely, assume \eqref{4.8.2}. Fix $t$ and $n$. Let's proceed by descending induction on $i$ to prove $\mathrm{H}^i(X_t, \mathcal{L}_t(n))$ vanishes. It vanishes for $i \gg 1$ by Serre's Theorem [\textbf{EGA III}$_1$, 2.2.2]. Suppose it vanishes for some $i \ge 2$. Then $\mathrm{R}^i f_{T*} (\mathcal{L}(n) \otimes_{\mathcal{O}_T} f_T^* \mathcal{N})_t$ vanishes for all quasi-coherent $\mathcal{N}$ on $T$, as just noted. So $\mathrm{R}^{i-1} f_{T*} (\mathcal{L}(n) \otimes_{\mathcal{O}_T} f_T^* \mathcal{N})_t$ is right exact in $\mathcal{N}$ owing to the long exact sequence of higher direct images. Therefore, by general principles, there is a natural isomorphism of functors
\[
\mathrm{R}^{i-1} f_{T*} (\mathcal{L}(n)) \otimes_{\mathcal{O}_T} \mathcal{N}_t \xrightarrow{\sim} \mathrm{R}^{i-1} f_{T*} (\mathcal{L}(n) \otimes_{\mathcal{O}_T} f_T^* \mathcal{N})_t.
\]
Since \eqref{4.8.2} holds, both source and target vanish. Taking $\mathcal{N} := k_t$ yields the vanishing of $\mathrm{H}^{i-1}(X_t, \mathcal{L}_t(n))$. Thus \eqref{4.8.2} implies \eqref{4.8.3}.

Finally, for any $t' \in T'$, any $i$, and any $n$, we have
\[
\mathrm{H}^i(X_{t'}, \mathcal{L}_{t'}(n)) = \mathrm{H}^i(X_{p(t')}, \mathcal{L}_{p(t')}(n)) \otimes_{k_t} k_{t'}
\]
because cohomology commutes with flat base change. So \eqref{4.8.3} remains valid after the base change $p : T' \to T$; whence, \eqref{4.8.2} does too. Thus the presheaf is well defined, and so $P^{\phi}_m$ is too.

Arguing much as we did for $P^{\phi} \times_P T$, we find, given a map $T \to P^{\phi}$, that, as $m$ varies, the products $P^{\phi}_m \times_{P^{\phi}} T$ form a nested sequence of open subschemes of $T$, whose union is $T$. In the argument, the key change is in proving openness. In place of [\textbf{EGA III}$_2$, 7.9.11], we use the following part of Serre's Theorem [\textbf{EGA III}$_1$, 2.2.2]: given a coherent sheaf $\mathcal{F}$ on a projective scheme over a Noetherian ring $A$, there are only finitely many $i \ge 1$ and $n \ge m$ such that $\mathrm{H}^i(\mathcal{F}(n))$ is nonzero, and all these nonzero $A$-modules are finitely generated. Hence, if there is a prime $\mathfrak{p}$ of $A$ such that $\mathrm{H}^i(\mathcal{F}(n))_{\mathfrak{p}} = 0$ for all $i \ge 1$ and $n \ge m$, then there is an $a \notin \mathfrak{p}$ such that $\mathrm{H}^i(\mathcal{F}(n))_a = 0$ for all $i \ge 1$ and $n \ge m$.

Therefore, again by [\textbf{EGA G} (0, 4.5.4), p. 103], it suffices to represent each $P^{\phi}_m$ by a quasi-projective $S$-scheme.

Fix $\phi$ and $m$. Set $\phi_0(n) := \phi(m+n)$. Then there is an isomorphism of functors $P^{\phi}_m \xrightarrow{\sim} P^{\phi_0}_0$, which is defined as follows. First, define an endomorphism $\varepsilon$ of $\Pic_{X/S}$ by sending an invertible sheaf $\mathcal{L}$ on an $X_T$ to $\mathcal{L}(m)$. Plainly $\varepsilon$ is an automorphism. So $\varepsilon$ induces an automorphism $\varepsilon^+$ of the associated sheaf $P$. Plainly $\varepsilon^+$ carries $P^{\phi}_m$ onto $P^{\phi_0}_0$. Thus it suffices to represent $P^{\phi_0}_0$ by a quasi-projective $S$-scheme.

The function $s \mapsto \chi(X_s, \mathcal{O}_{X_s}(n))$ is locally constant on $S$ by [\textbf{EGA III}$_2$, 7.9.11]. Hence we may assume it is constant by replacing $S$ by an open and closed subset. Set $\psi(n) := \chi(X_s, \mathcal{O}_{X_s}(n))$.

Consider the Abel map $\mathrm{A}_{X/S} : \Div_{X/S} \to P$. Note that $\Div_{X/S}$ is a scheme, in fact, an open subscheme of the Hilbert scheme $\Hilb_{X/S}$, by Theorem \textbf{3.7}. Form the product $P^{\phi_0}_0 \times_P \Div_{X/S}$. It is a scheme $Z$, in fact, an open subscheme of $\Div_{X/S}$, by what was proved above. Set $\theta(n) := \psi(n) - \phi_0(n)$. Plainly $Z$ lies in $\Hilb_{X/S}^{\theta}(n)$, which is projective over $S$. Hence $Z$ is quasi-projective.

Let's now prove the projection $\alpha : Z \to P^{\phi_0}_0$ is a surjection of \'{e}tale sheaves. In other words, given a $T$ and a $\lambda \in P^{\phi_0}_0(T)$, we have to find an \'{e}tale covering $T_1 \to T$ and a $\lambda_1 \in Z(T_1)$ such that $\alpha(\lambda_1) \in P^{\phi_0}_0(T_1)$ is equal to the image of $\lambda$.

Represent $\lambda$ by means of an \'{e}tale covering $p : T' \to T$ and an invertible sheaf $\mathcal{L}'$ on $X_{T'}$. Virtually by definition, the product $T' \times_{P^{\phi_0}_0} Z$ is equal to $\LinSys_{\mathcal{L}'/X_{T'}/T'}$. So, by Theorem \textbf{3.13}, this product is equal to $\mathbf{P}(\mathcal{Q})$ where $\mathcal{Q}$ is the $\mathcal{O}_{T'}$-module associated to $\mathcal{L}'$ as in Subsection \textbf{3.10}. Now, $m=0$, so $\mathrm{H}^1(X_t, \mathcal{L}_t) = 0$ for all $t \in T'$ owing to \eqref{4.8.3}. Since (v) implies (i) in Subsection \textbf{3.10} therefore $\mathcal{Q}$ is locally free.

Hence $\mathbf{P}(\mathcal{Q})$ is smooth over $T'$. So there exist an \'{e}tale covering $T_1 \to T'$ and a $T'$-map $T_1 \to \mathbf{P}(\mathcal{Q})$ by [\textbf{EGA IV}$_4$, 17.16.3 (ii)]. Then the composition $T_1 \to \mathbf{P}(\mathcal{Q}) \to Z \to P^{\phi_0}_0$ is equal to the composition $T_1 \to T' \to T \to P^{\phi_0}_0$. In other words, the map $T_1 \to Z$ is a $\lambda_1 \in Z(T_1)$ such that $\alpha(\lambda_1) \in P^{\phi_0}_0(T_1)$ is equal to the image of $\lambda$. Since the composition $T_1 \to T' \to T$ is an \'{e}tale covering, $\alpha$ is thus a surjection of \'{e}tale sheaves.

Plainly, the map $\alpha : Z \to P^{\phi_0}_0$ is defined by the invertible sheaf associated to the universal relative effective divisor on $X_Z/Z$. So taking $T := Z$ and $T' := T$ above, we conclude that the product $Z \times_{P^{\phi_0}_0} Z$ is a scheme and that the first projection is smooth and projective. Therefore, the theorem now results from the following general lemma. \hfill $\square$

\end{document}